\documentclass[../../main.tex]{subfiles}
\begin{document}

\begin{center}
  {\fontsize{48pt}{58pt}\selectfont \textbf{2012}\normalsize\textbf{年度}}
\end{center}

\begin{table}[htb]
  \centering
  \caption{出題テーマと難易度}
  \begin{tabular}{|c|p{9cm}|c|}
    \hline
    \textbf{問題番号} & \textbf{テーマ} & \textbf{難易度} \\
    \hline
    第1問 & 円板と半平面の共通部分を横切る線分の長さ & \\
    \hline
    第2問 & 三角形の部屋を移動する球の確率 & \\
    \hline
    第3問 & 領域の回転でできる二つの立体の体積 & \\
    \hline
    第4問 & 連続する自然数の積と累乗数 & \\
    \hline
    第5問 & 整数行列と面積1の平行四辺形 & \\
    \hline
    第6問 & 行列のトレースと回転行列 & \\
    \hline
  \end{tabular}
\end{table}

\begin{table}[htb]
  \centering
  \caption{解答}
  \renewcommand{\arraystretch}{1.4}
  \begin{tabular}{|c|c|p{9cm}|}
    \hline
    \textbf{問題番号} & \textbf{小問} & \textbf{答え} \\
    \hline
    第1問 & -- & $L$の最大値は$\dfrac{\sqrt{6}}{3}$，
    そのとき$\cos\theta=\dfrac{\sqrt{3}}{3}$ \\
    \hline
    第2問 & -- & $n$が奇数なら$0$，偶数なら
    $\dfrac{1}{3}\left\{1-\left(\dfrac{1}{2}\right)^{\frac{n}{2}}\right\}$ \\
    \hline
    第3問 & (1) & $V_1=\dfrac{11\pi}{480}$，
    $V_2=\dfrac{\pi}{192}(8\sqrt{2}-7)$ \\
    \cline{2-3}
    & (2) & $\dfrac{V_2}{V_1}=\dfrac{5}{22}(8\sqrt{2}-7)<1$ \\
    \hline
    第4問 & (1) & 連続する2自然数の積は$n$乗数でない（証明問題） \\
    \cline{2-3}
    & (2) & 連続する$n$自然数の積は$n$乗数でない（証明問題） \\
    \hline
    第5問 & (1)〜(3) & \\
    \hline
    第6問 & (1)・(2) & \\
    \hline
  \end{tabular}
\end{table}

\paragraph{試験形態}
6問．試験時間は要確認．

\end{document}
